\documentclass[../../../include/open-logic-section]{subfiles}

\begin{document}

\olfileid{sth}{replacement}{refproofs}
\olsection{পরিশিষ্ট: প্রতিস্থাপন-সংক্রান্ত ফল}

এই বিভাগে $\ZF$-এর মধ্যে প্রতিফলন প্রমাণ করব।
প্রতিফলন কোন অর্থে প্রতিস্থাপনের তুল্য,
তাও দেখাব। শেষে প্রতিফলন/প্রতিস্থাপনের
শক্তি সম্পর্কে একটি আকর্ষণীয় ফল প্রমাণ করব।
\emph{সতর্কতা: এই পাঠ্যবইয়ের সবচেয়ে
উন্নত গণিত এটাই।}

প্রথমে একটি সহায়ক ফল। সংক্ষেপের জন্য
এতে চলক বা বস্তুর অনুক্রম বোঝাতে
\emph{উপরি-রেখা} ব্যবহার করব। অর্থাৎ
``$\overline{a}_k$'' হলো
``$a_{k_1}$, \dots, $a_{k_n}$''-এর
সংক্ষিপ্ত রূপ; $n$ প্রসঙ্গ থেকে নির্ধারিত।

\begin{lem}\ollabel{lemreflection}
প্রত্যেক $1 \leq i \leq k$-এর জন্য
$\phi_i(\overline{v}_i, x)$ একটি সূত্র
হোক। তাহলে প্রত্যেক $\alpha$-র জন্য
এমন একটি $\beta > \alpha$ আছে যে
যেকোনো $\overline{a}_1,
\ldots, \overline{a}_k \in V_\beta$
ও প্রত্যেক $1 \leq i \leq k$-তে:
\[
	\exists x\phi_i(\overline{a}_i, x) \rightarrow (\exists x \in V_\beta) \phi_i(\overline{a}_i, x)
\]
\end{lem}

\begin{proof}
একটি পদ $\mu$ এইভাবে সংজ্ঞায়িত করি:
$\mu(\overline{a}_1, \ldots,
\overline{a}_k)$ হলো ক্ষুদ্রতম পর্যায়
$V$, যা প্রত্যেক $1 \leq i \leq k$-তে
নিচের শর্তবাক্যগুলি পূরণ করে:
% BN-SRC-442: উৎসের শর্তবাক্যে একটি বাড়তি বন্ধনী ছিল।
\[
\exists x\phi_i(\overline{a}_i, x) \rightarrow (\exists x \in V) \phi_i(\overline{a}_i, x)
\]
যেকোনো $\overline{a}_1, \ldots,
\overline{a}_k$-র জন্য
$\mu(\overline{a}_1, \ldots, \overline{a}_k)$
বিদ্যমান, তা সহজেই নিশ্চিত করা যায়।
এবার প্রতিস্থাপন ও আমাদের পুনরাবৃত্তির
উপপাদ্য ব্যবহার করে সংজ্ঞায়িত করি:
% BN-SRC-443: সংযোগের সূচক m, তাই সদস্য-পর্যায়ও S_m।
\begin{align*}
	S_0 & = V_{\alpha+1}\\
	S_{n+1} & = S_n \cup \bigcup
	\Setabs{\mu(\overline{a}_1, \ldots, \overline{a}_k)}
	{\overline{a}_1, \ldots, \overline{a}_k \in S_n} \\
	S &= \bigcup_{m < \omega} S_m.
\end{align*}
প্রত্যেক $S_n$, এবং তাই $S$-ও,
$V_\alpha$-র পরের একটি পর্যায়।
এবার $\overline{a}_1$, \dots,~$\overline{a}_k
\in S$ নিই; তাহলে কোনো $n < \omega$-তে
$\overline{a}_1$, \dots, $\overline{a}_k
\in S_n$। একটি $1 \leq i \leq k$
নিয়ে ধরি $\exists x
\phi_i(\overline{a}_i,x)$। সংজ্ঞা অনুযায়ী
$(\exists x \in \mu(\overline{a}_1,
\ldots, \overline{a}_k))\phi_i(\overline{a}_i, x)$;
তাই $(\exists x \in S_{n+1})\phi_i(\overline{a}_i, x)$,
ফলে $(\exists x \in S)\phi_i(\overline{a}_i, x)$।
অতএব এই $S$-ই চাওয়া $V_\beta$।
% BN-SRC-796 TARGET-CORRECTION-BEGIN
\emph{উৎস-সংশোধন-নোট।} প্রত্যেক স্থির finite tuple-তে সত্য
অস্তিত্ববাক্যগুলি থেকে সসীমসংখ্যক সাক্ষী নিয়ে তাদের স্তরাঙ্কের
উপরে এক পর্যায় পাই; মিথ্যা অস্তিত্ববাক্যের শর্ত শূন্যার্থে সত্য।
সূত্রভিত্তিক সর্বনিম্নতায় অনন্য ক্ষুদ্রতম পর্যায়, তাই $\mu$
মোট। Tuple-ক্ষেত্র সসীম Cartesian product, তার প্রতিচ্ছবি ও
স্বাভাবিক ধাপের পরিবার প্রতিস্থাপনে সেট। পর্যায়-পরিবারের সংযোগ
তার ক্রমসংখ্যা-সূচকগুলির supremum-এর পর্যায়; $S_n$ সঙ্গে
সংযোগেও তাই পর্যায় থাকে, শুরু $V_{\alpha+1}$-র নিচে নামে না।
সসীমসংখ্যক tuple-entry এই বাড়তে থাকা ধারায় এক সাধারণ
ধাপে থাকে। তারপর তার চাওয়া সাক্ষী পরের ধাপে যুক্ত হয়।
সব স্থির পরামিতি ভিতরে রাখতে প্রয়োজনে শুরুর সূচক বড় নিই।
% BN-SRC-796 TARGET-CORRECTION-END
\end{proof}
\noindent
এখন \olref[sth][replacement][ref]{reflectionschema}
বেশ সরাসরি প্রমাণ করা যায়:

\begin{proof}[প্রমাণ]
$\alpha$ নিই। সাধারণতার ক্ষতি ছাড়া
ধরতে পারি, $\phi$-র একমাত্র যৌক্তিক
ক্রিয়াগুলি $\exists$, $\lnot$ ও $\land$;
কারণ এগুলি দিয়ে সব সূত্র প্রকাশ করা যায়।
$\phi$-র উপসূত্রগুলিকে জটিলতা অনুসারে
$\psi_1, \ldots, \psi_k$ হিসেবে
তালিকাভুক্ত করি, যেখানে $\psi_k = \phi$।
\olref{lemreflection} অনুযায়ী এমন একটি
$\beta > \alpha$ আছে যে যেকোনো
$\overline{a}_i \in V_\beta$ ও
প্রত্যেক $1 \leq i \leq k$-তে:
\begin{align}\label{reflectionnicelybehaved}
	\exists x\psi_i(\overline{a}_i, x) \rightarrow
	(\exists x \in V_\beta) \psi_i(\overline{a}_i, x)\tag{*}
\end{align}
$\psi_i$-র জটিলতার উপর আবেশে দেখাব,
যেকোনো $\overline{a}_i \in V_\beta$-র জন্য
$\psi_i(\overline{a}_i) \leftrightarrow
\psi_i^{V_\beta}(\overline{a}_i)$।
$\psi_i$ পরমাণু সূত্র হলে কথাটি সরাসরি।
এই দ্বিশর্ত থেকেই বোঝা যায়,
$\psi_i$ এ ধর্মবিশিষ্ট উপসূত্রের নঞর্থকতা বা
সংযোজন হলে $\psi_i$-ও ধর্মটি পায়।
অতএব একমাত্র গুরুত্বপূর্ণ ক্ষেত্র
পরিমাণায়ন। $\overline{a}_i \in V_\beta$
নিই; তাহলে:
\begin{align*}
	(\exists x \psi_i(\overline{a}_i, x))^{V_\beta}
	&\text{ যদি ও কেবল যদি }
	(\exists x \in V_\beta)\psi_i^{V_\beta}(\overline{a}_i, x)
	&&\text{সংজ্ঞা অনুযায়ী}\\
	&\text{ যদি ও কেবল যদি }
	(\exists x \in V_\beta)\psi_i(\overline{a}_i,  x)
	&&\text{আবেশ-অনুমানে}\\
	&\text{ যদি ও কেবল যদি }
	\exists x \psi_i(\overline{a}_i, x)
	&&\text{\eqref{reflectionnicelybehaved} অনুযায়ী}
\end{align*}
এতে আবেশ শেষ; $\psi_k = \phi$ বলে
ফলটি অনুসৃত।
\end{proof}

$\ZF$-এ প্রতিফলন প্রমাণ করেছি।
আমাদের প্রমাণ মূলত \citet{Montague1961}-এর
পথ অনুসরণ করে। এবার $\Z$-এ
প্রতিফলন থেকে প্রতিস্থাপন প্রমাণ করতে
চাই। প্রমাণটি সরলীকরণসহ
\citet{Levy1960}-এর পথ অনুসরণ করে।

$\Z$-এ কাজ করছি বলে প্রতিফলনকে
উপরের ঠিক একই রূপে বলা যায় না:
উপরের রূপে ``$V_\alpha$'' সংকেত আছে,
কিন্তু $\Z$-এ সব সূচকে এর মোট মানের অস্তিত্ব প্রমাণিত নয়
(দেখুন \olref[sth][spine][zf]{sec})।
তাই প্রতিফলনের আপাতদৃষ্টিতে দুর্বলতর
একটি রূপ নিই:
% BN-SRC-444: উৎসে প্রতিস্থাপনের রূপ বলা; প্রদর্শিত নীতিটি দুর্বল প্রতিফলন।

\begin{defish}
\emph{দুর্বল প্রতিফলন।} যেকোনো সূত্র
$\phi$-র জন্য এমন একটি পরিযায়ী সেট
$S$ আছে যাতে $0$, $1$ এবং
$\phi$-র যেকোনো পরামিতি $S$-এর
!!{element}s; আর
$(\forall \overline{x} \in S)(\phi \liff
\phi^S)$।
\end{defish}

এটি ব্যবহার করে প্রতিস্থাপন প্রমাণের
জন্য প্রথমে \citet[উপপাদ্য ২-এর প্রথম অংশ]{Levy1960}
অনুসরণ করে দেখাই, একই সঙ্গে দুটি
সূত্র ``প্রতিফলিত'' করা যায়:

\begin{lem}[$\Z + \text{দুর্বল প্রতিফলন}$-এ]\ollabel{lem:reflect}
যেকোনো সূত্র $\psi, \chi$-র জন্য এমন
একটি পরিযায়ী সেট $S$ আছে যাতে
$0$ ও $1$ (এবং সূত্রগুলির যেকোনো
পরামিতি) $S$-এর !!{element}s;
আর $(\forall \overline{x} \in S)((\psi \liff
\psi^S) \land (\chi \liff \chi^S))$।
\end{lem}

\begin{proof}
$\phi$ হিসেবে $(z = 0 \land \psi) \lor (z = 1 \land \chi)$
সূত্রটি নিই।
% BN-SRC-797 TARGET-CORRECTION-BEGIN
\emph{উৎস-সংশোধন-নোট।} Tag-চল $z$ দুটি সূত্রের সমস্ত
মুক্ত বা বাঁধা চল থেকে নতুন বেছে নিই; প্রয়োজনে আগের বাঁধা
চলগুলি নামান্তর করি। তাই $z=0$ বা $z=1$ বসালে মূল
$\psi,\chi$-র চলের মান বদলায় না। ০ ও ১-এর পূর্ণ সূত্রে
সহায়ক বাঁধা চলগুলিও সংঘাতহীন। নিরপেক্ষতার এখানে ব্যবহৃত
সীমা $z\in S$; পরিযায়িতায় তার প্রত্যেক সদস্যও $S$-তে।
তাই এই দুই tag-সূত্রের পরিমাণায়ক সীমা বাইরের সদস্য হারায় না।
% BN-SRC-797 TARGET-CORRECTION-END

এখানে সংক্ষিপ্ত সংকেত ব্যবহার করেছি:
``$z = 0$''-কে ``$\forall t\, t \notin z$''
আর ``$z =1$''-কে ``$\forall s(s \in z
\liff \forall t\, t \notin s)$'' হিসেবে
পূর্ণ করে লেখা উচিত। কিন্তু $0, 1 \in S$
এবং $S$ পরিযায়ী বলে সূত্রগুলি
$S$-এর জন্য \emph{নিরপেক্ষ}:
পরিমাণায়কগুলি $S$-এ সীমাবদ্ধ করলেও
একই বস্তুতে এদের প্রয়োগ সত্য হয়।\footnote{আরও
আনুষ্ঠানিকভাবে, $\xi$-কে এই দুই
সূত্রের যেকোনো একটি ধরলে
$\xi(z) \liff \xi^S(z)$।}

দুর্বল প্রতিফলন অনুযায়ী উপযুক্ত
একটি $S$ পাই যাতে:
\begin{align*}
	(\forall z, \overline{x} \in S)(&\phi \liff \phi^S)\\
	\text{অর্থাৎ }(\forall z, \overline{x} \in S)(&((z = 0 \land \psi) \lor (z = 1 \land \chi)) \liff {}\\
	&\phantom{(}((z = 0 \land \psi) \lor (z = 1 \land \chi))^S)\\
	\text{অর্থাৎ }(\forall z, \overline{x} \in S)(&((z = 0 \land \psi) \lor (z = 1 \land \chi))\liff {}\\
	&\phantom{(}((z = 0 \land \psi^S) \lor (z = 1 \land \chi^S)))\\
	\text{অর্থাৎ }(\forall \overline{x} \in S)(&(\psi \liff \psi^S) \land (\chi \liff \chi^S))
\end{align*}
``$z = 0$'' ও ``$z=1$'' যে
$S$-এর জন্য নিরপেক্ষ, তার ফলে
দ্বিতীয় দাবি থেকে তৃতীয়টি আসে;
$0 \neq 1$ বলে চতুর্থ দাবি অনুসৃত।
\end{proof}\noindent
এখন \citet[উপপাদ্য ৬]{Levy1960}-এর
পদ্ধতি অনুসরণ ও সরলীকরণে
প্রতিস্থাপন পাওয়া যায়:

\begin{thm}[$\Z$ + দুর্বল প্রতিফলন-এ]\label{thm:replacement}
যেকোনো সূত্র $\phi(v,w)$ ও যেকোনো
$A$-র জন্য, যদি
$(\forall x \in A)\lexists![y][\phi(x,y)]$,
তবে
% BN-SRC-445: উৎসের set-builder শর্তে শেষ বন্ধনী অনুপস্থিত ছিল।
$\Setabs{y}{(\exists x \in A)\phi(x,y)}$
বিদ্যমান।
\end{thm}

\begin{proof}
এমন একটি $A$ নিই যাতে
$(\forall x \in A)\lexists![y][\phi(x,y)]$;
সূত্র দুটি সংজ্ঞায়িত করি:
\begin{align*}
	\psi &\text{ হলো } (\phi(x, z) \land A = A)\\
	\chi &\text{ হলো } \lexists[y][\phi(x, y)]
\end{align*}
\olref{lem:reflect} ব্যবহার করি।
$A$ যেহেতু $\psi$-র একটি পরামিতি,
তাই এমন একটি পরিযায়ী~$S$ আছে
যাতে $0, 1, A \in S$ (অন্য
পরামিতিগুলিও আছে), আর:
\[
	(\forall x,z \in S)((\psi \liff \psi^S) \land (\chi \liff \chi^S))
\]
ফলে বিশেষত:
\begin{align*}
	(\forall  x, z \in S)(&\phi(x, z) \liff \phi^S(x, z))\\
	(\forall x \in S)(&\exists y\phi(x, y) \liff (\exists y \in S)\phi^S(x, y))
\end{align*}
এই দুটি মেলাই; $A \in S$ এবং
$S$ পরিযায়ী বলে $A \subseteq S$-ও:
\begin{align*}
	(\forall x \in A)(&\exists y\phi(x, y) \liff (\exists y \in S)\phi(x, y))
\end{align*}
এখন $(\forall x \in A)(\lexists![y \in S])\phi(x, y)$,
কারণ $(\forall x \in A)\lexists![y][\phi(x, y)]$।
পৃথকীকরণ থেকে পাই
$\Setabs{y \in S}{(\exists x \in A) \phi(x, y)} = \Setabs{y}{(\exists
x \in A) \phi(x, y)}$।
\end{proof}

\end{document}
